\documentclass[addpoints]{exam}
\usepackage{amsmath}
\usepackage{amssymb}
\usepackage{geometry}
\geometry{margin=1in}

\begin{document}

\begin{center}
  {\Large \textbf{Security --- Final Exam}}\\[0.5em]
  {\large \textit{For each question, select \textbf{all} answers that apply.}}
\end{center}

\vspace{1em}

\begin{questions}
\question Which of the following are components of a URL?

\begin{checkboxes}
  \CorrectChoice Domain (e.g., \texttt{example.com})
  \choice Cookie
  \CorrectChoice Scheme (e.g., \texttt{http} or \texttt{https})
  \CorrectChoice Path (e.g., \texttt{/index.html})
  \CorrectChoice Port (e.g., \texttt{:8080})
\end{checkboxes}


\question The Same-Origin Policy (SOP) considers two pages to have the same origin when which of the following match?

\begin{checkboxes}
  \CorrectChoice Domain (e.g., \texttt{a.com} vs.\ \texttt{b.com})
  \choice Fragment identifier (\texttt{\#anchor})
  \CorrectChoice Scheme (e.g., \texttt{http} vs.\ \texttt{https})
  \choice Query string
  \CorrectChoice Port (e.g., \texttt{:80} vs.\ \texttt{:8080})
\end{checkboxes}


\question Which of the following are required properties of a cryptographic hash function?

\begin{checkboxes}
  \CorrectChoice Pre-image resistance
  \choice Key-dependent output
  \CorrectChoice Second pre-image resistance
  \choice Reversibility
  \CorrectChoice Collision resistance
\end{checkboxes}


\question Which of the following techniques can be used to bypass stack canaries?

\begin{checkboxes}
  \CorrectChoice Pointer subterfuge (overwriting a function pointer rather than the return address)
  \choice Using a NOP sled to slide over the canary
  \CorrectChoice Brute-forcing the canary value when the server forks and preserves the canary across crashes
  \choice Return-Oriented Programming (ROP) directly defeats a canary
  \choice Disabling DEP via \texttt{mprotect}
\end{checkboxes}


\question Which of the following are true about SQL injection attacks?

\begin{checkboxes}
  \choice SQL injection can only affect \texttt{SELECT} queries
  \CorrectChoice SQL injection exploits the mixing of executable code and untrusted data
  \choice HTTPS prevents SQL injection attacks
  \CorrectChoice An attacker can use SQL injection to bypass login authentication (e.g.,\ \texttt{OR '1'='1'})
  \CorrectChoice Prepared statements (parameterized queries) effectively prevent SQL injection
\end{checkboxes}


\question Which of the following are recognized types of stack canaries?

\begin{checkboxes}
  \choice Heap canaries
  \CorrectChoice Terminator canaries (contain bytes such as null or newline that terminate string operations)
  \choice NOP canaries
  \CorrectChoice Random canaries (a random value chosen at process start)
  \choice Return canaries
\end{checkboxes}


\question Which of the following are true about ARP (Address Resolution Protocol)?

\begin{checkboxes}
  \choice ARP messages are encrypted
  \CorrectChoice ARP maps IP addresses to MAC (hardware) addresses
  \CorrectChoice ARP spoofing can be used to perform man-in-the-middle attacks on a local network
  \choice ARP operates at the transport layer (Layer 4)
  \CorrectChoice ARP has no built-in authentication mechanism
\end{checkboxes}


\question Which of the following are defenses against DNS cache poisoning?

\begin{checkboxes}
  \CorrectChoice DNSSEC (cryptographically signing DNS records)
  \choice Increasing DNS TTL values
  \CorrectChoice DNS cookies (a cryptographic challenge-response mechanism)
  \CorrectChoice 0x20 encoding (randomizing the case of query names to add entropy)
  \CorrectChoice Randomizing the UDP source port used for DNS queries
\end{checkboxes}



\question Which of the following are established secure design principles?

\begin{checkboxes}
  \CorrectChoice Defense in depth
  \choice Security by obscurity
  \CorrectChoice Least privilege
  \CorrectChoice Complete mediation
  \CorrectChoice Privilege separation
\end{checkboxes}


\question Which of the following memory regions can be overwritten when exploiting a stack-based buffer overflow?

\begin{checkboxes}
  \CorrectChoice The return address
  \choice Heap metadata
  \CorrectChoice Local variables on the stack
  \CorrectChoice Function arguments on the stack
  \CorrectChoice The saved frame pointer (EBP)
\end{checkboxes}


\question Which of the following are types of Denial-of-Service (DoS) or DDoS attacks?

\begin{checkboxes}
  \CorrectChoice HTTP flood
  \choice SQL injection
  \CorrectChoice SYN flood
  \CorrectChoice Logic-based DoS (exploiting bugs to exhaust server resources)
  \CorrectChoice Reflection and amplification attacks
\end{checkboxes}


\question Which of the following are x86 general-purpose registers?

\begin{checkboxes}
  \CorrectChoice EAX
  \choice EIP
  \CorrectChoice ECX
  \CorrectChoice EDX
  \CorrectChoice EBX
\end{checkboxes}


\question Which of the following are true about TLS 1.3 compared to TLS 1.2?

\begin{checkboxes}
  \choice TLS 1.3 uses RSA for key exchange
  \CorrectChoice TLS 1.3 encrypts more of the handshake, hiding certificate information from passive observers
  \choice TLS 1.3 requires clients to present a certificate for authentication
  \CorrectChoice TLS 1.3 includes downgrade protection to prevent forcing the use of older, weaker protocol versions
  \choice TLS 1.3 is backwards compatible with SSL v2
\end{checkboxes}


\question Which of the following are true about Cross-Site Scripting (XSS)?

\begin{checkboxes}
  \CorrectChoice XSS allows an attacker to execute malicious scripts in a victim's browser
  \choice XSS directly exploits the server's database
  \CorrectChoice In reflected XSS, the malicious script is reflected off the server in a response
  \choice XSS requires the victim to be logged into the targeted site
  \CorrectChoice Content Security Policy (CSP) can help mitigate XSS attacks
\end{checkboxes}


\question Which of the following are true about Return-Oriented Programming (ROP)?

\begin{checkboxes}
  \CorrectChoice ROP chains together small code sequences (``gadgets'') that each end in a \texttt{ret} instruction
  \choice ROP requires injecting new shellcode into the process
  \CorrectChoice ROP can bypass DEP/W\^{}X protections because it reuses existing executable code
  \choice ROP is only possible on 64-bit systems
  \CorrectChoice ROP is Turing-complete
\end{checkboxes}



\question Which of the following are true about DNS cache poisoning?

\begin{checkboxes}
  \CorrectChoice An attacker can inject forged DNS responses into a recursive resolver's cache
  \choice DNS cache poisoning requires physical access to the DNS server
  \CorrectChoice Bailiwick checking limits what records a name server is allowed to include in a response
  \choice HTTPS prevents DNS cache poisoning
  \CorrectChoice The Kaminsky attack targets an entire zone by flooding a resolver with forged responses for random subdomains
\end{checkboxes}


\question Which of the following are true about Subresource Integrity (SRI)?

\begin{checkboxes}
  \choice SRI is applied automatically by all browsers without any developer action
  \CorrectChoice SRI uses cryptographic hashes embedded in HTML tags to verify third-party resource content
  \choice SRI prevents all forms of Cross-Site Scripting
  \CorrectChoice SRI can prevent malicious code injection when a CDN is compromised
  \choice SRI encrypts third-party resources in transit
\end{checkboxes}



\question Which of the following encryption approaches or modes are considered insecure?

\begin{checkboxes}
  \choice Cipher Block Chaining (CBC) mode with a random IV
  \CorrectChoice Electronic Codebook (ECB) mode
  \CorrectChoice Caesar cipher (easily broken by frequency analysis)
  \choice AES-256
  \CorrectChoice Reusing a one-time pad key for more than one message
\end{checkboxes}



\question Which of the following are true about Diffie-Hellman (DH) key exchange?

\begin{checkboxes}
  \CorrectChoice DH allows two parties to establish a shared secret over an insecure public channel
  \choice DH requires a pre-shared secret to initiate the exchange
  \CorrectChoice The security of DH relies on the hardness of the discrete logarithm problem
  \choice DH directly encrypts messages between the two parties
  \CorrectChoice DH is vulnerable to an active man-in-the-middle attack when used without authentication
\end{checkboxes}


\question Which of the following are true about BGP (Border Gateway Protocol)?

\begin{checkboxes}
  \CorrectChoice BGP is used for routing between different Autonomous Systems on the Internet
  \CorrectChoice BGP hijacking can redirect Internet traffic through an attacker-controlled network
  \choice BGP provides end-to-end encryption of routed traffic
  \CorrectChoice BGP has no built-in mechanism to authorize route announcements
  \choice BGP is only used within a single organization's internal network
\end{checkboxes}



\question Which of the following are true about UNIX file permissions?

\begin{checkboxes}
  \CorrectChoice The \texttt{setuid} bit allows a program to run with the file owner's effective privileges
  \choice All files default to world-readable unless explicitly restricted
  \CorrectChoice Each file has three permission sets: owner, group, and others
  \choice Only the root user can set the execute bit on a file
  \CorrectChoice Permissions include read (r), write (w), and execute (x) bits
\end{checkboxes}


\question Which of the following are techniques used for web tracking?

\begin{checkboxes}
  \CorrectChoice Canvas fingerprinting
  \choice DNS over HTTPS
  \CorrectChoice Cookie syncing between ad networks
  \CorrectChoice Tracking cookies
  \CorrectChoice Browser fingerprinting
\end{checkboxes}


\question Which of the following are true about HTTP?

\begin{checkboxes}
  \CorrectChoice HTTP status code 302 indicates a redirect
  \choice HTTP encrypts all traffic by default
  \CorrectChoice HTTP is a stateless protocol
  \CorrectChoice HTTP status code 404 indicates that a resource was not found
  \CorrectChoice HTTP supports multiple request methods including GET and POST
\end{checkboxes}



\question Which of the following are valid defenses against Cross-Site Request Forgery (CSRF)?

\begin{checkboxes}
  \CorrectChoice CSRF tokens (unpredictable tokens embedded in forms)
  \choice Content Security Policy (CSP)
  \CorrectChoice Using the \texttt{SameSite} cookie attribute
  \choice SQL prepared statements
  \CorrectChoice Validating the \texttt{Referer} or \texttt{Origin} request header
\end{checkboxes}


\question Which of the following are valid DNS record types?

\begin{checkboxes}
  \CorrectChoice A (maps a hostname to an IPv4 address)
  \choice HTTP
  \CorrectChoice MX (specifies a mail server for a domain)
  \CorrectChoice NS (specifies the authoritative name server for a domain)
  \CorrectChoice AAAA (maps a hostname to an IPv6 address)
\end{checkboxes}


\question Which of the following are true about the Bleichenbacher attack?

\begin{checkboxes}
  \choice It targets the discrete logarithm problem
  \CorrectChoice It is a padding oracle attack
  \CorrectChoice It exploits PKCS\#1 v1.5 RSA padding
  \choice It attacks HMAC constructions
  \CorrectChoice It is an adaptive chosen-ciphertext attack
\end{checkboxes}



\question Which of the following are components of the CIA triad?

\begin{checkboxes}
  \CorrectChoice Confidentiality
  \choice Authentication
  \CorrectChoice Integrity
  \choice Non-repudiation
  \CorrectChoice Availability
\end{checkboxes}


\question Which of the following are true about HTTP cookies?

\begin{checkboxes}
  \CorrectChoice Cookies are set by the server using the \texttt{Set-Cookie} response header
  \choice The browser validates the authenticity of cookie contents
  \CorrectChoice Persistent cookies have a specified expiration date
  \choice Cookies are encrypted by default
  \CorrectChoice Session cookies are deleted when the browser session ends
\end{checkboxes}


\question Which of the following are true about the Rowhammer attack?

\begin{checkboxes}
  \CorrectChoice Rowhammer exploits bit-flip vulnerabilities in DRAM
  \choice Rowhammer requires kernel-level code execution
  \CorrectChoice Error Correcting Code (ECC) memory can help mitigate Rowhammer
  \choice Rowhammer targets the CPU cache rather than main memory
  \CorrectChoice The attack works by repeatedly accessing rows adjacent to the target row
\end{checkboxes}


\question Which of the following are true about Certificate Authorities (CAs) and the PKI?

\begin{checkboxes}
  \CorrectChoice Root CA certificates are hardcoded into browsers and operating systems
  \choice CAs can decrypt TLS traffic to any site for which they issued a certificate
  \CorrectChoice A CA signs a server's public key (in a certificate) to vouch for its identity
  \CorrectChoice A misissued certificate for a domain can enable a man-in-the-middle attack against that domain
  \CorrectChoice Intermediate CAs can issue certificates on behalf of a root CA
\end{checkboxes}


\question Which of the following are true about Message Authentication Codes (MACs)?

\begin{checkboxes}
  \CorrectChoice HMAC is a widely used MAC construction based on a cryptographic hash function
  \choice MACs provide confidentiality
  \CorrectChoice MACs require a shared secret key
  \CorrectChoice MACs provide message authentication
  \CorrectChoice MACs provide message integrity verification
\end{checkboxes}


\question Which of the following are true about firewalls and network filtering?

\begin{checkboxes}
  \CorrectChoice Packet-filtering firewalls inspect IP and transport-layer headers (e.g., source/destination IP and port)
  \choice Firewalls can decrypt HTTPS traffic without a man-in-the-middle certificate
  \CorrectChoice Proxy-based firewalls can inspect and filter application-layer content
  \CorrectChoice Network Address Translation (NAT) implicitly blocks unsolicited inbound connections
  \choice A firewall alone protects against all application-layer attacks
\end{checkboxes}



\question Which of the following security properties does TLS provide?

\begin{checkboxes}
  \CorrectChoice Confidentiality (encryption of data in transit)
  \choice Non-repudiation of all transactions
  \CorrectChoice Integrity (detection of tampering with data in transit)
  \choice Anonymity of the client to the server
  \CorrectChoice Authentication of the server via its certificate
\end{checkboxes}


\question Which of the following describe the adversarial mindset in security?

\begin{checkboxes}
  \choice Optimizing the system for performance over correctness
  \CorrectChoice Questioning false assumptions made by defenders
  \CorrectChoice Identifying and targeting the weakest link in a system
  \choice Trusting user-supplied input by default
  \CorrectChoice Thinking outside the box to find unexpected attack paths
\end{checkboxes}



\question Which of the following are true about TCP?

\begin{checkboxes}
  \CorrectChoice TCP provides reliable, ordered data delivery
  \choice TCP provides built-in encryption
  \CorrectChoice TCP uses a three-way handshake (\texttt{SYN}, \texttt{SYN-ACK}, \texttt{ACK}) to establish connections
  \choice TCP prevents IP address spoofing
  \CorrectChoice TCP uses sequence numbers to order segments and detect packet loss
\end{checkboxes}


\question Which of the following are true about unsafe C functions and memory-safe languages?

\begin{checkboxes}
  \CorrectChoice \texttt{strcpy} does not check destination buffer bounds and is commonly associated with buffer overflows
  \choice \texttt{strncpy} is fully safe and always null-terminates the destination
  \CorrectChoice \texttt{gets} reads input without a length limit and is considered unsafe
  \choice Buffer overflows are only possible in interpreted languages
  \CorrectChoice Rust and Go are examples of memory-safe languages that help prevent buffer overflows
\end{checkboxes}



\question Which of the following are true about the Mirai botnet?

\begin{checkboxes}
  \CorrectChoice Mirai primarily compromised IoT devices (e.g., IP cameras and routers)
  \choice Mirai primarily targeted desktop Windows PCs
  \CorrectChoice Mirai used a command-and-control (C2) server architecture to direct attacks
  \choice Mirai exploited previously unknown (zero-day) software vulnerabilities
  \CorrectChoice Mirai exploited devices using default or weak credentials
\end{checkboxes}


\question Which of the following \texttt{printf} format specifiers can be exploited in a format string attack?

\begin{checkboxes}
  \CorrectChoice \texttt{\%x} (can leak values off the stack)
  \choice \texttt{\%f} (reads a floating-point value; not exploitable for writes)
  \CorrectChoice \texttt{\%n} (writes the number of bytes printed to a memory address)
  \CorrectChoice \texttt{\%s} (can be used to read arbitrary memory as a string)
  \choice \texttt{\%p} is the only format specifier that can leak stack values
\end{checkboxes}


\question Which of the following are mitigations for format string and integer overflow vulnerabilities?

\begin{checkboxes}
  \CorrectChoice Using strongly typed languages that enforce type safety
  \choice Stack canaries
  \CorrectChoice Static analysis to detect unsafe format string usage
  \CorrectChoice Runtime integer overflow checking
  \choice Address Space Layout Randomization (ASLR)
\end{checkboxes}



\question Which of the following are true about virtual memory and hardware privilege levels on x86?

\begin{checkboxes}
  \CorrectChoice The Memory Management Unit (MMU) translates virtual addresses to physical addresses using page tables
  \CorrectChoice Ring 0 is the most privileged level (kernel mode) on Intel hardware
  \choice Virtual memory requires no hardware support
  \CorrectChoice Each process has its own virtual address space, providing isolation
  \CorrectChoice Hypervisors can run at Ring $-1$, more privileged than the OS kernel
\end{checkboxes}



\question Which of the following are examples of side channels?

\begin{checkboxes}
  \CorrectChoice Power consumption
  \choice File permissions
  \CorrectChoice Execution time
  \CorrectChoice Memory access patterns
  \CorrectChoice Electromagnetic emissions
\end{checkboxes}


\question Which of the following are true about Meltdown and Spectre?

\begin{checkboxes}
  \CorrectChoice Kernel Page Table Isolation (KPTI) is a mitigation for Meltdown
  \CorrectChoice Spectre exploits speculative execution triggered by branch prediction
  \choice Both vulnerabilities require the attacker to be running code in the OS kernel
  \CorrectChoice Both attacks use a cache side channel to exfiltrate speculatively accessed data
  \CorrectChoice Meltdown exploits out-of-order execution and can allow user-space to read kernel memory
\end{checkboxes}



\question Which of the following are valid responses to an identified risk?

\begin{checkboxes}
  \CorrectChoice Mitigate the risk
  \choice Escalate the risk
  \CorrectChoice Avoid the risk
  \CorrectChoice Accept the risk
  \CorrectChoice Transfer the risk
\end{checkboxes}


\question Which of the following are true about integer overflow vulnerabilities in C?

\begin{checkboxes}
  \choice Using \texttt{int} instead of \texttt{unsigned int} prevents integer overflow
  \CorrectChoice Unsigned integer overflow causes wrap-around behavior
  \CorrectChoice Truncation can occur when a larger integer type is assigned to a smaller one
  \choice Integer overflows are automatically detected and raised as exceptions in C
  \CorrectChoice Signed-to-unsigned type conversions can produce unexpected results
\end{checkboxes}


\question Which of the following are true about RSA?

\begin{checkboxes}
  \CorrectChoice The RSA private key is the pair $(d, N)$
  \choice RSA decryption uses the public key
  \CorrectChoice Textbook RSA (without padding) is malleable and insecure in practice
  \CorrectChoice The RSA public key is the pair $(e, N)$
  \CorrectChoice RSA key generation involves selecting two large prime numbers $p$ and $q$, and computing $N = p \cdot q$
\end{checkboxes}


\end{questions}

\end{document}
